% !TeX program = xelatex
\documentclass[11pt,letterpaper]{article}
\usepackage[margin=1in]{geometry}
\usepackage{iftex}
\ifPDFTeX
  \usepackage[utf8]{inputenc}
  \usepackage[T1]{fontenc}
  \usepackage{helvet}
  \renewcommand{\familydefault}{\sfdefault}
  \DeclareUnicodeCharacter{00A7}{\S}
  \DeclareUnicodeCharacter{00B7}{\textperiodcentered}
  \DeclareUnicodeCharacter{2013}{--}
  \DeclareUnicodeCharacter{2014}{---}
  \DeclareUnicodeCharacter{201C}{``}
  \DeclareUnicodeCharacter{201D}{''}
  \DeclareUnicodeCharacter{2193}{\(\downarrow\)}
\else
  \usepackage{fontspec}
\fi
\usepackage[spanish,es-nodecimaldot]{babel}
\usepackage{xcolor}
\usepackage{graphicx}
\usepackage{array}
\usepackage{booktabs}
\usepackage{longtable}
\usepackage{enumitem}
\usepackage{fancyhdr}
\usepackage{titlesec}
\usepackage{tcolorbox}
\usepackage{hyperref}
\usepackage{fvextra}
\graphicspath{{./}{imagenes/}{images/}{figures/}}
\ifPDFTeX
\else
  \IfFontExistsTF{Arial}{
    \setmainfont{Arial}
    \setsansfont{Arial}
  }{
    \setmainfont{TeX Gyre Heros}
    \setsansfont{TeX Gyre Heros}
  }
\fi
\definecolor{FabricBlack}{HTML}{0A0A0A}
\definecolor{FabricPanel}{HTML}{161616}
\definecolor{FabricChampagne}{HTML}{C9A96E}
\definecolor{FabricText}{HTML}{252525}
\hypersetup{colorlinks=true, linkcolor=FabricBlack, urlcolor=FabricChampagne}
\pagestyle{fancy}
\fancyhf{}
\lhead{\textsf{FABRIC Oracle Critical Engineering}}
\rhead{\textsf{An\'alisis y Dise\~no}}
\cfoot{\thepage}
\titleformat{\section}{\Large\bfseries\sffamily\color{FabricBlack}}{\thesection}{0.8em}{}[\titlerule]
\titleformat{\subsection}{\large\bfseries\sffamily\color{FabricBlack}}{\thesubsection}{0.8em}{}
\titleformat{\subsubsection}{\normalsize\bfseries\sffamily\color{FabricBlack}}{\thesubsubsection}{0.8em}{}
\setlist[enumerate]{leftmargin=*, itemsep=0.25em, topsep=0.25em}
\setlength{\parindent}{0pt}
\setlength{\parskip}{0.65em}
\newcommand{\diagramimage}[3]{
\begin{figure}[htbp]
\centering
\IfFileExists{#1}{
  \includegraphics[width=\textwidth, height=0.88\textheight, keepaspectratio]{#1}
}{
  \IfFileExists{imagenes/#1}{
    \includegraphics[width=\textwidth, height=0.88\textheight, keepaspectratio]{imagenes/#1}
  }{
    \fbox{\parbox[c][0.28\textheight][c]{0.88\textwidth}{
    \centering
    \textbf{Imagen pendiente}\\[0.4em]
    Sube el archivo \texttt{#1} a la raíz del proyecto o a la carpeta \texttt{imagenes/}.
    }}
  }
}
\caption{#2}
\label{#3}
\end{figure}
}

\begin{document}
\begin{titlepage}
\pagecolor{FabricBlack}\color{white}
\vspace*{2cm}
{\Huge\bfseries\sffamily FABRIC\par}
\vspace{0.3cm}
{\Large Oracle Critical Engineering\par}
\vfill
{\huge\bfseries Documento de An\'alisis y Dise\~no\par}
\vspace{0.35cm}
{\Large FABRIC / fabricsoft.com.mx\par}
\vspace{1cm}
{\large Versi\'on t\'ecnica revisada -- Motor IA + Panel de Administraci\'on en V1\par}
{\large Mayo 2026\par}
\vfill
{\color{FabricChampagne}\rule{\textwidth}{1.2pt}}\par
\vspace{0.4cm}
{\large Equipo A -- Edi + Gerardo -- Supervisi\'on Julio \'Alvarez\par}
\end{titlepage}
\nopagecolor\color{FabricText}
\tableofcontents
\newpage

Demo intermedia: jueves 21 de mayo · Demo final: lunes 25 de mayo · Standup: 11:00 AM

\subsection{Índice ejecutivo}

\begin{longtable}{p{0.46\linewidth}p{0.46\linewidth}}
\toprule
\textbf{Sección} & \textbf{Contenido} \\
\midrule
Sección 1 Contexto del negocio y dirección estratégica & Misión, objetivo general, objetivos específicos, requerimientos, reglas, metas y alcance. \\
Sección 2 Análisis de mercado y perfil del usuario & Mercado objetivo, perfiles CFO/CTO y casos de uso principales. \\
Sección 3 Arquitectura de la información y navegación & Sitemap, navegación persistente y criterios de orientación del usuario (prueba del tronco). \\
Sección 4 Especificaciones de UI/UX & Sistema visual, tokens de diseño, restricciones visuales y lineamientos de formularios. \\
Sección 5 Maquetado y prototipado visual & Entregables de maquetado, criterios de aprobación visual y estados de interfaz. \\
Sección 6 Capa de presentación y arquitectura técnica & Stack tecnológico, estructura de componentes, renderizado y glosario de términos. \\
Sección 7 Motor IA — especificación técnica & Arquitectura del agente Oracle, orquestador de modelos, RAG, Kill Switch y High-Intent Routing. \\
Sección 8 Panel de administración & Acceso seguro con Clerk, roles de usuario, gestión de leads y métricas en vivo. \\
Sección 9 Integración y consumo de datos & Proveedores de servicios, estimación de costos, modelo MongoDB y flujos de lead magnets. \\
Sección 10 Infraestructura, rendimiento y despliegue & Estrategia Lighthouse 95+, seguridad perimetral, hosting y optimización SEO técnico. \\
\bottomrule
\end{longtable}

\section{Contexto del negocio y dirección estratégica}

Este documento establece los lineamientos estratégicos y técnicos para la arquitectura del ecosistema digital de \textbf{FABRIC} (\texttt{fabricsoft.com.mx}) en su Versión 1 (V1), enfocada en el desarrollo e integración del Agente Oracle (Motor de Inteligencia Artificial) y el Centro de Mando Administrativo interno.

\subsection{Misión y visión}

\subsubsection{Misión}
Acompañar y liderar la remediación de proyectos Oracle críticos garantizando estabilidad operativa durante el primer ciclo transaccional en producción, erradicando el enfoque transitorio del tradicional ``go-live''.

\subsubsection{Visión}
Posicionar a FABRIC como la firma de ingeniería crítica Oracle líder en México y Latinoamérica, reconocida por CFOs y CTOs por su capacidad para mitigar riesgos, asegurar la continuidad operativa y recuperar implementaciones ERP inestables.

\subsection{Posicionamiento estratégico y doctrina operativa}

\subsubsection{Matriz de Posicionamiento Maestro}
El modelo de negocio de FABRIC se fundamenta en un posicionamiento claro que redefine la relación consultora-cliente en proyectos complejos:
\begin{description}
\item[Categoría Madre:] Oracle Critical Engineering.
\item[Promesa Contractual:] \emph{``No entregamos en go-live. Entregamos cuando tu primer ciclo crítico opera en producción.''}
\item[Enemigo Común:] El integrador tradicional cuya responsabilidad expira con el ``go-live'' del sistema, dejando al cliente desamparado frente a las inconsistencias contables y operativas subsiguientes.
\item[Tagline de Marca:] Oracle Cloud, entregado en producción.
\item[Público Objetivo:] Directores de Finanzas (CFO) y Directores de Tecnología (CTO/CIO) de corporaciones con facturación de USD \$50M a \$500M+ que enfrentan proyectos fallidos de Oracle Fusion, EBS, JDE o PeopleSoft.
\end{description}

\subsubsection{Fórmula del ROI Técnico}
El valor entregado por FABRIC se cuantifica a través del control presupuestal y la erradicación del reprocesamiento manual post go-live:
\[ \text{ROI Técnico} = \frac{\Delta \text{Horas de Cierre Contable} + \text{Costo de Reportes Manuales Eliminados}}{\text{Inversión en Remediación Fixed-Price}} \]
Donde la inversión es fija por fase, garantizando previsibilidad absoluta al CFO.

\subsection{Objetivo general del ecosistema digital}
Establecer la plataforma web oficial de FABRIC como un filtro de autoridad técnica y gatekeeper operativo. El diseño rompe con los esquemas de ventas B2B tradicionales; en su lugar, invierte la dinámica de poder, requiriendo que sea el prospecto quien califique y justifique la necesidad de los servicios de ingeniería de la firma.

\subsection{Objetivos específicos del proyecto}
\begin{enumerate}
\item \textbf{Validación Autónoma de Autoridad:} Demostrar la profundidad técnica de la firma en tiempo real a través de la consola interactiva del Agente Oracle.
\item \textbf{Filtrado Server-Side Pre-Contacto:} Implementar un motor de admisión estricto que redirija prospectos no calificados de forma silenciosa para optimizar el ancho de banda del equipo comercial.
\item \textbf{Conversión Orientada a Datos:} Registrar diagnósticos de dolor operativo (\emph{Rescue Assessment}), simulaciones financieras de TCO e interacciones de alta intención en base de datos.
\item \textbf{Evidencia Verificable y Transparencia:} Exponer métricas de rendimiento real de la firma y descargas controladas de papers de investigación técnica en un portal dedicado.
\item \textbf{Excelencia Técnica Demostrable:} Garantizar un rendimiento técnico intachable con puntuaciones Lighthouse de 95+ y métricas Web Vitals óptimas (LCP < 1.5s).
\item \textbf{Autonomía Operativa:} Proveer un panel administrativo para controlar slots de asesoría, métricas públicas y umbrales de capacidad sin requerir redespliegues de código.
\end{enumerate}

\subsection{Ciclo de vida del servicio (FABRIC Lifecycle)}
El modelo de prestación de servicios de ingeniería de FABRIC no concluye con la implementación o salida a producción de un sistema ERP. Se estructura en cinco fases bien delimitadas, las cuales determinan el comportamiento del software y de la base de datos (MongoDB) en la captura y catalogación de leads:
\begin{enumerate}
\item \textbf{01 · Diagnose (Diagnóstico):} Evaluación asíncrona a través del wizard web (\emph{Rescue Assessment}) y estimación financiera de desvíos presupuestales y operativos.
\item \textbf{02 · Architect (Arquitectura):} Diseño detallado de la estrategia de remediación, interfaces contables, optimizaciones y la recopilación de datos para RAG.
\item \textbf{03 · Deploy (Despliegue):} Construcción y puesta en producción del software remedial y las automatizaciones requeridas.
\item \textbf{04 · Stabilize (Estabilización):} Acompañamiento directo en producción durante la ejecución del primer ciclo contable y operativo crítico, garantizando la inmutabilidad y exactitud de los saldos sin reprocesos manuales.
\item \textbf{05 · Optimize (Optimización):} Ajustes continuos de rendimiento, análisis de datos acumulados y automatización incremental mediante agentes IA.
\end{enumerate}

\subsection{Requerimientos técnicos}

\subsubsection{Performance obligatorio}

\begin{longtable}{p{0.31\linewidth}p{0.31\linewidth}p{0.31\linewidth}}
\toprule
\textbf{Métrica} & \textbf{Umbral} & \textbf{Herramienta} \\
\midrule
Lighthouse Performance & 95+ & Lighthouse CI + Vercel Analytics \\
Lighthouse Accessibility & 95+ & Lighthouse CI \\
Lighthouse SEO & 100 & Lighthouse CI \\
LCP — Largest Contentful Paint & < 1.5s & Web Vitals \\
CLS — Cumulative Layout Shift & < 0.05 & Web Vitals \\
TTFB — Time to First Byte & < 200ms & Vercel Edge \\
FID — First Input Delay & < 50ms & Web Vitals \\
\bottomrule
\end{longtable}

\subsubsection{Requerimientos arquitectónicos}

\begin{longtable}{p{0.46\linewidth}p{0.46\linewidth}}
\toprule
\textbf{Requerimiento} & \textbf{Implementación} \\
\midrule
Arquitectura desacoplada & Separar presentación en Next.js, contenido en Sanity, datos operativos en MongoDB y Motor IA como módulo independiente. \\
Streaming asíncrono de baja latencia & El agente IA responde mediante Vercel AI SDK sin bloquear la interfaz. El usuario visualiza la respuesta conforme se genera. \\
Validación del lado servidor & La validación de email, revenue, cargo, honeypot y reglas de admisión se ejecuta en Server Actions o API routes. Zod valida antes de escribir en base de datos. \\
Datos editables sin deploy & Métricas, slots y estado de capacidad se consumen desde Sanity o MongoDB mediante ISR y panel admin. \\
Alta disponibilidad y Edge & Vercel Edge Network, Cloudflare DNS, Static Generation por defecto, ISR en páginas con contenido dinámico y Edge Functions para formularios críticos. \\
Panel admin protegido & Rutas /admin/* detrás de middleware Next.js + Clerk, con restricción de dominio corporativo y roles. \\
\bottomrule
\end{longtable}

\subsubsection{Requerimientos de código}

\begin{enumerate}
\item \textbf{Tipado Estricto:} Uso obligatorio de TypeScript en modo estricto (\texttt{strict: true}), prohibiendo el uso de tipos implícitos o explícitos de tipo \texttt{any} para asegurar la robustez del código.
\item \textbf{Validación Servidor:} Implementación obligatoria de esquemas de validación con \texttt{Zod} en el lado del servidor para todas las peticiones a API Routes y Server Actions, validando los datos antes de cualquier escritura en la base de datos.
\item \textbf{Documentación y Modularidad:} Cada componente de presentación debe incluir un bloque de comentarios inicial detallando su propósito, dependencias y comportamiento esperado, facilitando la mantenibilidad futura del ecosistema.
\end{enumerate}

\subsection{Reglas de negocio y lógica de filtrado}

\subsubsection{Arquitectura de entornos y flujos de precalificación}
El ecosistema digital opera bajo tres entornos independientes con responsabilidades claras y flujos de datos asíncronos:
\begin{longtable}{p{0.46\linewidth}p{0.46\linewidth}}
\toprule
\textbf{Entorno} & \textbf{Descripción y acceso} \\
\midrule
Sitio público (\texttt{fabricsoft.com.mx}) & Interfaz de presentación, terminal del Agente Oracle, wizard de diagnóstico y calculadores de TCO. No requiere autenticación. \\
Centro de mando (\texttt{/admin}) & Consola privada protegida por Clerk para gestión de leads calificados, slots de asesoría, métricas operativas y logs de auditoría. \\
CMS (\texttt{Sanity Studio}) & Interfaz editorial para la administración descentralizada de casos de éxito, papers técnicos, variables del comparador y slots de agenda. \\
\bottomrule
\end{longtable}

\subsubsection{Doctrina FABRIC V4.0 — Compromisos operativos y contractuales}
La operación de la firma y el comportamiento del software están regidos por cinco principios inmutables que se integran en la lógica de negocio del sistema:
\begin{enumerate}
\item \textbf{Principio de Entrega en Primer Ciclo Crítico (vs Go-Live):} El hito de cierre de proyecto no se define por la salida a producción (go-live), sino por el funcionamiento estable del primer ciclo transaccional crítico (ej. cierre contable mensual). El sistema de tickets del panel administrativo no marcará un caso como ``Cerrado'' hasta que este ciclo se complete satisfactoriamente.
\item \textbf{Cero Juniors Facturables:} La asignación de recursos en la firma se limita a ingenieros y consultores con un mínimo de 8 años de experiencia real verificable en sistemas Oracle. Los perfiles en el CMS y las asignaciones en bases de datos deben cumplir obligatoriamente esta restricción de seniority.
\item \textbf{Presupuesto Fixed-Price por Fase:} Todos los alcances se cotizan bajo un esquema de precio fijo. Si se presentan desviaciones en los tiempos de entrega imputables a la firma, el sistema de facturación y el panel administrativo bloquean la generación de cargos adicionales o semanas de soporte extraordinarias.
\item \textbf{Cero Procesamiento o Reporte Manual Post-Go-Live:} FABRIC garantiza que todos los reportes y procesos contables comprometidos queden automatizados dentro del ERP. Si por limitaciones de diseño subsisten conciliaciones manuales por causas imputables a la firma, FABRIC las resolverá sin coste alguno.
\item \textbf{Transición con Documentación Viva:} Cada entrega de proyecto incluye un tablero operativo de KPIs, runbooks específicos de incidentes, configuraciones del sistema y matriz de roles. Estos entregables deben cargarse en el repositorio seguro asociado al cliente antes de la firma del acta de transición.
\end{enumerate}

\subsubsection{Motor de admisión y filtros server-side (Zod + API Validation)}
El formulario de aplicación (\texttt{/aplicar}) y el wizard de precalificación implementan reglas automáticas de filtrado del lado del servidor para garantizar que solo los prospectos que califiquen como clientes ideales puedan agendar una sesión:
\begin{enumerate}
\item \textbf{Restricción de Dominios Genéricos:} Bloqueo estricto a nivel de esquema Zod para correos que contengan dominios personales o públicos (Gmail, Yahoo, Hotmail, etc.), requiriendo obligatoriamente un correo con dominio corporativo.
\item \textbf{Filtro de Facturación Corporativa:} Los prospectos con ingresos declarados menores a USD \$50M anuales son clasificados en la base de datos de MongoDB con el estado \texttt{waitlist} de forma automatizada y transparente.
\item \textbf{Validación de Intención y Plazo:} Se requiere que el prospecto seleccione un plazo de iniciativa Oracle activo menor a 6 meses. Plazos mayores o consultorías indefinidas son marcados como baja prioridad.
\item \textbf{Patrocinio Ejecutivo Requerido:} La validación del formulario exige que se declare tanto el patrocinio del Director de Finanzas (CFO) como del Director de Tecnología (CTO/CIO). La ausencia de alguno de estos roles en la iniciativa suspende el agendamiento directo.
\end{enumerate}

\subsubsection{Doctrina de filtrado — Criterios de exclusión}
El motor de precalificación descarta y redirige a la lista de espera general a cualquier contacto que cumpla con los siguientes criterios de exclusión de la firma (Doctrina de No Alineación):
\begin{itemize}
\item Pequeñas y medianas empresas (PYMES) que carezcan de un departamento formal de cumplimiento regulatorio y auditoría.
\item Iniciativas orientadas a soporte preventivo simple, actualización de parches sin remediación de procesos, o administración básica de infraestructura.
\item Solicitudes de subcontratación de personal (staff augmentation o body shopping) por tarifa horaria.
\item Proyectos que no cuenten con el patrocinio activo del CFO del negocio.
\end{itemize}

\subsubsection{Motor IA (Agente Oracle) — Reglas operativas}
\begin{enumerate}
\item \textbf{Delimitación del Dominio Técnico:} El Agente Oracle limitará estrictamente sus respuestas a temas relacionados con Oracle Fusion Cloud, Oracle EBS, JDE, PeopleSoft, OCI y casos de estudio o metodologías de FABRIC.
\item \textbf{Interruptor de Emergencia (Kill Switch):} El sistema detendrá y bloqueará en tiempo real consultas que estén fuera del alcance temático, intentos de inyección de prompts (jailbreaks) o solicitudes inadecuadas, sugiriendo el contacto con un ingeniero.
\item \textbf{Orquestación y Fallback Multi-modelo:} Para asegurar la disponibilidad continua, la lógica del orquestador intentará procesar las consultas prioritariamente en Claude Sonnet. Ante una latencia superior a 3 segundos o error de servicio, reanudará la petición mediante GPT-4o, y posteriormente con Grok como segundo respaldo.
\item \textbf{Enrutamiento de Alta Intención (High-Intent Routing):} El motor analizará semánticamente la conversación en busca de señales de urgencia crítica (e.g., caídas del sistema en producción, cierres contables bloqueados, multas). Al detectarlas, activará de inmediato un botón destacado para agendar Office Hours.
\item \textbf{Preparación de RAG:} La arquitectura del agente queda configurada para recuperar contexto relevante a partir de embeddings vectoriales del corpus oficial de Oracle y los manuales operativos de FABRIC.
\item \textbf{Privacidad de Datos Corporativos:} Los datos recolectados en diagnósticos y conversaciones se cifran tanto en tránsito como en reposo, garantizando que no serán compartidos ni utilizados para entrenar modelos de IA externos.
\end{enumerate}

\subsubsection{Panel admin — Reglas de acceso y control de capacidad}
\begin{enumerate}
\item \textbf{Autenticación Corporativa Estricta:} Acceso exclusivo al panel privado mediante Clerk, restringido a cuentas con el dominio corporativo \texttt{@fabricsoft.com.mx} y con verificación de dos factores (2FA) obligatoria.
\item \textbf{Control de Acceso Basado en Roles (RBAC):} Diferenciación de perfiles de usuario: \textit{Super Admin} (permisos de edición total sobre métricas, capacidad y accesos) y \textit{Admin} (permisos limitados a lectura y gestión comercial de leads).
\item \textbf{Trazabilidad e Inmutabilidad (Audit Trail):} Registro inalterable en la base de datos de todas las operaciones críticas realizadas por los administradores, almacenando identificador, tipo de acción, fecha y snapshot del estado anterior y posterior.
\item \textbf{Embudo de Escasez y Capacidad Límite:} El número máximo de proyectos activos de la firma se configura en el panel. Si se alcanza dicho límite, el sitio web desactiva dinámicamente el widget de Calendly para agendar Office Hours y habilita el llamado de acción hacia la lista de espera.
\item \textbf{Slots de Asesoría Rolling:} Se limitan los slots mensuales de consultoría estratégica a un máximo de 4 slots rolling para mantener el enfoque exclusivo de los ingenieros principales de la firma.
\end{enumerate}

\subsection{Metas e indicadores de éxito}

\begin{longtable}{p{0.46\linewidth}p{0.46\linewidth}}
\toprule
\textbf{Meta} & \textbf{Indicador} \\
\midrule
Precalificación efectiva & Leads con dominio corporativo y datos mínimos de empresa, cargo, revenue e iniciativa. \\
Validación técnica autónoma & Eventos de analítica: ai\_query\_completed, diagnostic\_completed y paper\_downloaded. \\
Precalificación de presupuesto & El lead llega con ERP actual, revenue declarado, tipo de iniciativa y resultado del diagnóstico. \\
Gestión de demanda & Office Hours ocupadas antes de la mitad del mes y activación controlada de Wait List cuando corresponda. \\
Excelencia técnica & CI aprobado antes de demos; sin scroll horizontal en mobile; LCP <1.5s verificado. \\
Panel admin operativo & Demo final con visualización de leads, aprobación/rechazo, edición de métricas y control de capacidad. \\
\bottomrule
\end{longtable}

\subsection{Alcance}

\subsubsection{Entregables contemplados para el 25 de mayo}

\begin{enumerate}
\item Home completa con 16 secciones S01-S16.
\item 15 páginas internas activas y 9 páginas Próximamente con captura de email.
\item Cinco lead magnets funcionales: Rescue Assessment, ERP Comparator, Wait List, Office Hours y Papers.
\item Motor IA Oracle con RAG preparado, Kill Switch, fallback y High-Intent Routing.
\item Panel Admin con Clerk para gestión de leads, métricas y capacidad.
\item Sanity Studio con contenido inicial: casos, papers, FSOs, métricas y slots.
\item Staging en equipoa.fabricsoft.com.mx operativo antes de la demo intermedia.
\end{enumerate}

\subsubsection{Fuera de alcance para V1}

\begin{enumerate}
\item Ejecución de scripts o cambios sobre infraestructura real de clientes.
\item Pasarela de pago, e-commerce o CRM avanzado.
\item API pública para partners o integraciones externas.
\item Aplicación móvil nativa.
\item Precios cerrados; las condiciones comerciales se negocian fuera del sitio.
\end{enumerate}

\section{Análisis de mercado y perfil del usuario}

\subsection{Segmentación y mercado objetivo}
El mercado objetivo de FABRIC se limita a organizaciones que operan sistemas complejos Oracle (Fusion Cloud, EBS, JDE, PeopleSoft) en México y Latinoamérica, con facturación anual superior a USD \$50M. La firma concentra su autoridad técnica y comunicación en tres sectores industriales exclusivos:

\begin{longtable}{p{0.25\linewidth}p{0.35\linewidth}p{0.32\linewidth}}
\toprule
\textbf{Vertical Focal} & \textbf{Características de Operación} & \textbf{Clientes de Referencia} \\
\midrule
Servicios Financieros y Fintech & Procesamiento masivo de transacciones, conciliaciones bancarias automatizadas de alta velocidad, estricto cumplimiento ante la CNBV y auditoría transaccional. & ABC Capital, Aplazo, Copayment \\
\midrule
Inmobiliario y Centros Comerciales & Complejidad en la facturación y cobranza de miles de contratos de arrendamiento, adendas variables, cálculo de mantenimiento y consolidación de ingresos multi-sitio. & APE Plazas (Caso ancla 2026) \\
\midrule
Logística, Distribución y Transporte & Gestión integrada de cadena de suministro (SCM), inventarios en tiempo real, facturación de fletes, interfaces de transporte (TMS) y optimización de rutas. & ALMEX, AJE, ENEX, Arcos Dorados \\
\bottomrule
\end{longtable}

\subsection{Caso de Éxito Ancla: APE Plazas (Estabilización en Producción)}
El proyecto de remediación para \textbf{APE Plazas} en abril de 2026 representa la materialización de la doctrina operativa de FABRIC. 

\subsubsection{Contexto y Problemática Operativa}
Tras la salida a producción (go-live) inicial realizada por otro integrador, APE Plazas enfrentó un bloqueo operativo crítico: las interfaces de facturación de arrendamientos con el motor de contratos de Oracle Fusion fallaron, dejando más de 12,000 locales comerciales sin facturación y forzando al departamento de finanzas a realizar conciliaciones en hojas de cálculo externas. El go-live técnico se había celebrado, pero el negocio estaba paralizado.

\subsubsection{Ejecución de Ingeniería de Remediación}
FABRIC asumió la dirección técnica en la fase de remediación. El equipo de ingenieros principales realizó las siguientes intervenciones de software y arquitectura:
\begin{itemize}
\item \textbf{Reingeniería del Motor de Reglas Contables (SLA):} Rediseño de las reglas de contabilización de contratos en Oracle Fusion Financials para soportar adendas variables.
\item \textbf{Optimización de Integración de API:} Refactorización de las Serverless Functions que conectaban el sistema operativo de rentas con las APIs REST de Oracle ERP Cloud, reduciendo el índice de fallas de integración de un 34\% a menos del 0.05\%.
\item \textbf{Cierre Contable Estabilizado:} Acompañamiento en sitio y soporte proactivo para procesar el primer ciclo crítico correspondiente al mes de \textbf{Abril de 2026}.
\end{itemize}
El hito final del proyecto se alcanzó el 30 de abril de 2026, logrando el primer cierre contable mensual completamente automatizado dentro de Oracle Fusion, sin reprocesamientos manuales y con trazabilidad total para auditoría.

\subsection{Buyer persona y perfiles de decisión}

\subsubsection{Persona A — CFO (Comprador Económico)}
\textbf{Perfil:} Director Financiero o Vicepresidente de Finanzas. Su principal métrica es la predictibilidad de costos, el cumplimiento de auditorías y la reducción del costo total de propiedad (TCO) del ERP. Su mayor frustración son los desvíos presupuestales y los procesos manuales persistentes que maquillan implementaciones incompletas.

\textbf{Comportamiento en la Plataforma:} El CFO utiliza la plataforma para verificar la seriedad financiera y legal de FABRIC. Analiza la doctrina contractual de precio fijo (Fixed-Price) y descarga papers relacionados con el ROI y los riesgos de implementación de Oracle Fusion para sustentar la inversión ante el consejo.

\subsubsection{Persona B — CTO / CIO (Comprador Técnico)}
\textbf{Perfil:} Director de Tecnología o Sistemas. Responsable de la estabilidad de la infraestructura, de la disponibilidad del ERP y de la integración de plataformas corporativas. Exige rigor de ingeniería, documentación de sistemas, y su mayor dolor es la asignación de consultores junior que no resuelven incidentes complejos.

\textbf{Comportamiento en la Plataforma:} El CTO interactúa directamente con la consola de comandos del Agente Oracle para auditar la competencia técnica de la firma en tiempo real. Analiza las capas de arquitectura técnica (\emph{FABRIC OS}) y descarga notas de investigación y patrones de falla técnica de Oracle en LATAM.

\subsection{Casos de uso core}

\begin{longtable}{p{0.31\linewidth}p{0.31\linewidth}p{0.31\linewidth}}
\toprule
\textbf{Caso de uso} & \textbf{Actor} & \textbf{Módulos / resultado} \\
\midrule
CU1 — Diagnóstico técnico autónomo & CTO / CIO & S04 y Agente IA Oracle. Resultado: validación técnica antes del contacto humano. \\
CU2 — Intercambio de valor & CFO & ERP TCO, Cloud y Rescue Assessment. Resultado: lead con datos financieros y contexto de iniciativa. \\
CU3 — Filtro y agendamiento & Prospecto C-Level de alta intención & Office Hours y Wait List. Resultado: agenda visible solo si pasa criterios mínimos. \\
CU4 — Centro de mando privado & Super Admin y equipo FABRIC & /admin con Clerk. Resultado: leads, métricas y capacidad operativa gestionables sin tocar código. \\
CU5 — Consumo de investigación técnica & CTO / CFO & Descarga controlada de Papers. Resultado: validación del correo corporativo, registro en la BD y envío asíncrono vía Resend. \\
\bottomrule
\end{longtable}

\subsection{Diagrama de casos de uso}

El sistema público funciona como un filtro activo de prospectos y como una demostración técnica de autoridad. El panel privado concentra la operación interna: revisión de leads, métricas, capacidad, agenda y trazabilidad.

\diagramimage{diagrama-casos-uso.jpeg}{Diagrama de casos de uso de FABRIC.}{fig:casos-uso}

\subsubsection{Descripción ejecutiva de actores}

\begin{longtable}{p{0.22\linewidth}p{0.34\linewidth}p{0.34\linewidth}}
\toprule
\textbf{Actor} & \textbf{Interés principal} & \textbf{Resultado esperado} \\
\midrule
Visitante corporativo & Entender si FABRIC resuelve su problema y si su empresa califica. & Navegación clara, propuesta de valor visible y filtro no invasivo. \\
CFO & Validar impacto financiero, riesgo, ROI y seriedad comercial. & Comparadores, casos, papers y señales de control presupuestal. \\
CTO / CIO & Verificar profundidad técnica y capacidad real de ejecución Oracle. & Agente IA, arquitectura, papers y diagnóstico técnico autónomo. \\
Equipo FABRIC Admin & Operar demanda, leads, capacidad y métricas sin redeploy. & Panel privado con acciones auditables y datos centralizados. \\
Servicios externos & Autenticación, IA, notificaciones y sincronización comercial. & Integración segura mediante API routes y Server Actions. \\
\bottomrule
\end{longtable}

\section{Arquitectura de la información y navegación}

Referencia visual: fabric\_wireframes\_s01\_s16.html, entregable con las 16 secciones anotadas para desarrollo.

\subsection{Sitemap — resumen ejecutivo}

\begin{longtable}{p{0.31\linewidth}p{0.31\linewidth}p{0.31\linewidth}}
\toprule
\textbf{Tipo} & \textbf{Cantidad} & \textbf{Estado V1} \\
\midrule
Home — 16 secciones S01-S16 & 1 página & Funcional completa \\
Páginas internas activas & 15 páginas & Incluidas en wireframes \\
Páginas Próximamente & 9 páginas & UI premium + email capture \\
Panel Admin /admin/* & 1 interfaz privada & Clerk + dashboard Next.js \\
Sanity Studio /studio & 1 interfaz editorial & Contenido editable \\
\bottomrule
\end{longtable}

\subsection{Navegación persistente}

\begin{enumerate}
\item Logo \textbf{FABRIC} renderizado con la tipografía corporativa \emph{GT Sectra Display} con enlace dinámico a la página principal.
\item Enlaces de navegación principal: Doctrina, Casos, Industrias, FABRIC OS, Diagnóstico y Aplicar.
\item Botón de Acción Principal (CTA) ``Aplicar'' en color champagne con efecto de subrayado expansivo de 300 ms en hover.
\item Punto de control mobile (breakpoint a 968 px) con menú colapsable tipo overlay a pantalla completa.
\item Atributo de accesibilidad \emph{Skip-to-content} como primer elemento focable del DOM.
\end{enumerate}

\subsection{Prueba del tronco}

En cada página, el usuario debe poder responder sin esfuerzo: dónde está, qué es FABRIC, qué problema resuelve, por qué debería confiar y cuál es el siguiente paso. Estas respuestas deben aparecer en el primer viewport o quedar claramente insinuadas desde la estructura de navegación.

\section{Especificaciones de UI/UX}

El diseño debe comunicar especialización, control y criterio técnico. La interfaz evita clichés visuales B2B y prioriza claridad, sobriedad y rendimiento.

\subsection{Tokens de diseño}

\subsubsection{Paleta de colores}

\begin{longtable}{p{0.31\linewidth}p{0.31\linewidth}p{0.31\linewidth}}
\toprule
\textbf{Variable CSS} & \textbf{Valor} & \textbf{Uso} \\
\midrule
--bg-base & \#0A0A0A & Fondo base de la aplicación. \\
--bg-panel & \#161616 & Fondo de tarjetas, menús y navbar con scroll. \\
--bg-elevated & \#1A1A1A & Elementos flotantes o elevados. \\
--border & \#2A2A2A & Bordes y divisiones secundarias. \\
--border-strong & \#353535 & Bordes de elementos activos o enfocados. \\
--text-primary & \#F5F5F5 & Títulos, etiquetas y cuerpo principal. \\
--text-secondary & \#8A8A8A & Textos complementarios y metadatos. \\
--text-tertiary & \#5A5A5A & Placeholders y textos inactivos. \\
--accent & \#C9A96E & Color primario de acento (Champagne). \\
--accent-2 & \#A07845 & Color de acento secundario (Bronce). \\
--danger & \#B85450 & Alertas del sistema y errores de validación. \\
\bottomrule
\end{longtable}

\subsubsection{Especificación tipográfica}
El sistema tipográfico está optimizado para garantizar legibilidad, elegancia y alto rendimiento mediante fuentes cargadas localmente con la directiva \texttt{next/font}:
\begin{longtable}{p{0.25\linewidth}p{0.45\linewidth}p{0.22\linewidth}}
\toprule
\textbf{Familia Tipográfica} & \textbf{Uso en la Interfaz} & \textbf{Rango de Tamaño} \\
\midrule
\emph{GT Sectra Display} & Titulares mayores de secciones y logotipos (serif elegante). & 72px -- 120px \\
\emph{GT Sectra Text} & Subtítulos, citas y destacados tipográficos. & 32px -- 48px \\
\emph{Inter} o \emph{Söhne} & Cuerpo de texto general, párrafos y descripciones. & 17px -- 18px \\
\emph{JetBrains Mono} & Consolas, bloques de código, tablas de datos y logs. & 14px -- 15px \\
\emph{Söhne Mono Caps} & Etiquetas, badges y metadatos con tracking +5\%. & 12px \\
\bottomrule
\end{longtable}

\subsubsection{Restricciones visuales}
El sistema prohíbe explícitamente elementos comunes en el diseño web SaaS clásico para mantener un perfil formal de ingeniería:
\begin{enumerate}
\item \textbf{Cero Colores Saturados:} Prohibido el uso de azul SaaS (\texttt{\#3B82F6}), verde fluorescente o rojo neón.
\item \textbf{Cero Gradientes Coloridos:} Las transiciones de fondo deben ser sólidas y sobrias.
\item \textbf{Cero Estilos de Moda Visual:} Prohibido el uso de \emph{glassmorphism} (fondos translúcidos con desenfoque) o \emph{neumorphism} (relieves de sombra blanda).
\item \textbf{Cero Fotografías de Stock:} Solo se permite iconografía técnica precisa y diagramas de arquitectura reales.
\item \textbf{Bordes Angulados:} No se permiten bordes redondeados con un radio superior a 4 px en botones, paneles o campos de entrada de datos.
\end{enumerate}

\subsection{Microinteracciones y política de movimiento}
El sitio sigue una doctrina de ``quietud visual'' para transmitir seriedad y evitar distracciones operativas. Las animaciones están severamente restringidas a transiciones de interfaz indispensables:
\begin{enumerate}
\item \textbf{Revelación de Contenido (Scroll Reveal):} Transición de opacidad ($0 \to 1$) y desplazamiento vertical ($24\,\text{px} \to 0$), con una duración máxima de 200 ms y aceleración de salida (\emph{ease-out}).
\item \textbf{Efecto Typewriter en Terminal:} Simulación de comandos en la terminal de IA con velocidad de 30 ms por carácter y cursor parpadeante mediante CSS.
\item \textbf{Subrayado Expansivo (CTAs):} Línea de acento champagne que se expande desde el centro hacia los extremos en 300 ms al hacer hover.
\item \textbf{Prohibiciones de Movimiento:} Queda prohibido el uso de efectos de desplazamiento de fondo (\emph{parallax}), control forzado de scroll (\emph{scroll-jacking}), y animaciones de celebración (como lluvia de confeti) en los envíos de formularios.
\end{enumerate}

\subsection{Formularios de alta conversión}

\begin{longtable}{p{0.46\linewidth}p{0.46\linewidth}}
\toprule
\textbf{Principio} & \textbf{Implementación} \\
\midrule
Una pregunta por pantalla & Wizard secuencial; no mostrar formularios largos completos. \\
Progreso visible & Indicador “Paso N de N” en JetBrains Mono y barra champagne animada. \\
Validación campo a campo & Errores inline, específicos y visibles en el momento. \\
Email corporativo & Validación cliente y servidor contra dominios genéricos. \\
Anti-spam multicapa & Honeypot, tiempo mínimo, Turnstile invisible y rate limit por IP. \\
Confirmación post-submit & Correo transaccional con Resend en menos de 30 segundos cuando aplique. \\
\bottomrule
\end{longtable}

\section{Maquetado y prototipado visual}
Esta seccion queda reservada para incorporar el maquetado final del sitio y sus pantallas clave. El objetivo es conectar los lineamientos de UI/UX con decisiones visuales verificables antes de pasar a implementacion.

\subsection{Entregables de maquetado}
\begin{enumerate}
\item Home completa con las secciones S01--S16 en vista desktop, tablet y mobile.
\item Flujos principales: diagnostico, comparador ERP, aplicacion a Wait List, Office Hours y descarga de papers.
\item Estados de interfaz: loading, error, exito, rechazo silencioso, capacidad llena y navegacion mobile.
\item Panel Admin con dashboard, tabla de leads, edicion de metricas y control de capacidad.
\end{enumerate}

\subsection{Criterios de aprobacion visual}
\begin{enumerate}
\item Coherencia con tokens definidos: fondo oscuro, acento champagne, bordes sobrios y tipografia premium.
\item Lectura clara en primer viewport: autoridad, problema que resuelve FABRIC y accion siguiente.
\item Diseno responsive sin scroll horizontal, sin solapamientos y con jerarquia clara en pantallas pequenas.
\item Prototipo navegable apto para revision ejecutiva y handoff tecnico.
\end{enumerate}

\section{Arquitectura técnica — capa de presentación}

La arquitectura propuesta separa presentación, autenticación, ejecución serverless, servicios de IA, seguridad perimetral y persistencia operativa. Esta división permite iterar el sitio público sin comprometer la seguridad del panel administrativo ni la integridad de los datos capturados.

\diagramimage{arquitectura-stack.jpeg}{Arquitectura visual del stack tecnológico: Next.js, React, TypeScript, Tailwind CSS, AI SDK, MongoDB Atlas, Cloudflare, Clerk y Vercel.}{fig:arquitectura-stack}

\subsection{Stack seleccionado}

\begin{longtable}{p{0.31\linewidth}p{0.31\linewidth}p{0.31\linewidth}}
\toprule
\textbf{Tecnología} & \textbf{Versión} & \textbf{Justificación} \\
\midrule
Next.js & 15 · App Router & SSG, ISR, API Routes y Server Actions con buena base para rendimiento y SEO. \\
React & 19 & Server Components reducen JavaScript enviado al cliente. \\
TypeScript & 5.x strict & Mejora auditabilidad, mantenimiento y calidad de código. \\
Tailwind CSS & 3.x & Permite implementar tokens del brief en tailwind.config.ts con purge automático. \\
Framer Motion & 11.x & Se limita a microinteracciones aprobadas y tree-shakeables. \\
Vercel AI SDK & 4.x & Streaming del agente IA mediante useChat. \\
Clerk & 5.x & Autenticación para panel admin, middleware y roles. \\
next/font & Built-in & Carga controlada de fuentes y reducción de FOUT. \\
next/image & Built-in & WebP automático, lazy loading y prevención de CLS. \\
\bottomrule
\end{longtable}

\subsection{Arquitectura de componentes}

\begin{Verbatim}[fontsize=\small, breaklines=true]
fabricsoft/
├── app/
│   ├── page.tsx                 # Home — Secciones S01 a S16
│   ├── layout.tsx               # Root layout, metadatos y fuentes
│   ├── globals.css              # Variables CSS del sistema de diseño
│   ├── casos/                   # Páginas de Casos de Éxito
│   ├── doctrina/                # Páginas de Doctrina Operativa
│   ├── industrias/              # Páginas de Verticales Focales
│   ├── fabric-os/               # Páginas de Metodología FABRIC OS
│   ├── diagnostico/
│   │   └── page.tsx             # Formulario de Rescue Assessment (Wizard)
│   ├── comparator/
│   │   └── erp/page.tsx         # Comparador TCO de ERP actual vs Oracle
│   ├── aplicar/
│   │   └── page.tsx             # Formulario de aplicación a Wait List
│   ├── office-hours/
│   │   └── page.tsx             # Agenda de slots calificados vía Calendly
│   ├── admin/                   # Dashboard privado protegido por Clerk
│   └── api/                     # Endpoints, Server Actions y Agente IA
├── components/
│   ├── layout/                  # Nav, Footer, MobileMenu y SkipToContent
│   ├── sections/                # Componentes individuales de S01 a S16
│   ├── admin/                   # Componentes del panel (LeadsTable, etc.)
│   ├── forms/                   # Wizards y validaciones de formularios
│   └── ui/                      # Elementos base (Button, ProgressBar, etc.)
├── lib/
│   ├── sanity.ts                # Conector cliente de Sanity CMS
│   ├── mongodb.ts               # Cliente base de datos MongoDB Atlas
│   ├── resend.ts                # Configuración de correos con Resend
│   ├── hubspot.ts               # Integración con CRM HubSpot
│   ├── ai.ts                    # Lógica y fallbacks del orquestador de IA
│   ├── validators.ts            # Esquemas de validación Zod
│   ├── tco-calculator.ts        # Algoritmo de cálculo de ahorro TCO
│   └── analytics.ts             # Eventos personalizados para Plausible
└── middleware.ts                # Control de accesos y protección de /admin/*
\end{Verbatim}

\subsection{Estrategia de render por página}

\begin{longtable}{p{0.31\linewidth}p{0.31\linewidth}p{0.31\linewidth}}
\toprule
\textbf{Página} & \textbf{Estrategia} & \textbf{Razón} \\
\midrule
/ & Static Generation & No requiere datos críticos en primer render. \\
/casos/ape-plazas · /transparencia & ISR revalidate 3600 s & Contenido de Sanity con cambios ocasionales. \\
/diagnostico · /comparator/erp & Client-side & Wizard interactivo; SEO secundario. \\
/aplicar · /office-hours & Static + Server Action & Formulario estático con submit seguro del lado servidor. \\
/admin/* & Server-side + Clerk middleware & Datos operativos y autenticación obligatoria. \\
Páginas Próximamente & Static Generation & Contenido fijo y máximo rendimiento. \\
\bottomrule
\end{longtable}

\subsection{Explicación simple de tecnologías y herramientas}

El siguiente cuadro detalla las herramientas seleccionadas para el proyecto, traducidas a su función práctica y justificación de negocio para facilitar el entendimiento de la dirección general:

\begin{longtable}{p{0.20\linewidth}p{0.38\linewidth}p{0.36\linewidth}}
\toprule
\textbf{Tecnología} & \textbf{¿Qué es en palabras simples?} & \textbf{¿Por qué se usa en FABRIC?} \\
\midrule
Next.js & Framework principal para construir sitios web interactivos y veloces. & Combina excelente SEO (visible para Google), velocidad de carga rápida y ejecución segura en el servidor. \\
React & Biblioteca de programación para construir la interfaz visual. & Permite crear componentes de pantalla interactivos y reutilizables, como formularios, botones y widgets. \\
TypeScript & Versión estructurada de JavaScript que añade reglas estrictas. & Evita errores de programación comunes en tiempo de desarrollo y simplifica el mantenimiento a largo plazo. \\
Tailwind CSS & Sistema utilitario para el diseño de estilos de la interfaz. & Asegura consistencia visual de la marca (colores, espaciados y tipografías) con una alta velocidad de desarrollo. \\
Framer Motion & Biblioteca especializada en animaciones interactivas. & Añade micro-interacciones sobrias y elegantes que mejoran la experiencia de usuario sin saturar el navegador. \\
Vercel & Plataforma de alojamiento y despliegue continuo en la nube. & Permite actualizaciones instantáneas sin interrupción del servicio y excelente velocidad de respuesta global. \\
Cloudflare & Escudo de seguridad perimetral y administración de dominios (DNS). & Protege el sitio contra ataques maliciosos y abusos de bots, y ofrece un CAPTCHA invisible (Turnstile). \\
Sanity CMS & Administrador de contenidos dinámico y autohospedado. & Permite al equipo de FABRIC editar casos de éxito, artículos y métricas visibles en tiempo real sin requerir cambios de código. \\
MongoDB Atlas & Base de datos flexible en la nube. & Almacena de forma segura leads calificados, resultados de simulaciones, historiales de chat IA y logs de auditoría. \\
Clerk & Proveedor externo de autenticación y gestión de usuarios. & Protege el acceso al panel administrativo limitando la entrada solo a correos con dominio corporativo mediante 2FA. \\
Zod & Validador estricto de estructuras de datos en el servidor. & Garantiza que la información enviada en los formularios cumpla con los formatos requeridos antes de guardarla. \\
Resend & Servicio optimizado de envío de correos electrónicos transaccionales. & Envía de forma inmediata los PDFs de comparativas de TCO, notificaciones de contacto y papers técnicos. \\
HubSpot CRM & Plataforma de gestión de relaciones con clientes. & Centraliza los leads calificados que genera el sitio web para el seguimiento comercial del equipo de FABRIC. \\
Calendly & Sistema automatizado para agendar reuniones. & Permite a prospectos pre-calificados reservar de manera directa slots de asesoría técnica de alta prioridad. \\
Plausible & Sistema de analítica web simplificado enfocado en privacidad. & Mide conversiones y clics críticos sin almacenar cookies invasivas ni datos personales de los usuarios. \\
Sentry & Monitor de errores de software en tiempo real. & Envía alertas automáticas al equipo de ingeniería si ocurre algún fallo en el backend o frontend para resolverlo de inmediato. \\
\bottomrule
\end{longtable}

\subsection{Glosario ejecutivo de conceptos técnicos}

Esta guía define términos técnicos avanzados que aparecen a lo largo de este documento, traduciendo su impacto técnico a valor para el negocio:

\begin{longtable}{p{0.24\linewidth}p{0.70\linewidth}}
\toprule
\textbf{Concepto} & \textbf{Explicación para dirección y negocio} \\
\midrule
Arquitectura desacoplada & División del sistema en módulos independientes (web, base de datos, CMS, IA). Si una parte falla o requiere mantenimiento, el resto del sitio continúa operativo. \\
Server-side (Lado servidor) & Procesamiento de datos y lógica que ocurre en servidores protegidos. Es indispensable para resguardar reglas de negocio y validar datos sensibles. \\
Client-side (Lado cliente) & Ejecución que ocurre directamente en el navegador del visitante. Se utiliza para animaciones, transiciones suaves y validaciones visuales rápidas. \\
ISR (Incremental Static Regeneration) & Permite actualizar la información de ciertas secciones (como cifras del Rescue Counter) en segundos sin necesidad de reconstruir o compilar todo el sitio. \\
SSG (Static Site Generation) & Pre-compilación estática de páginas web. Genera pantallas ultra-rápidas que cargan al instante y benefician directamente el SEO del sitio. \\
API Route (Ruta de API) & Endpoint en el servidor que recibe peticiones del navegador para realizar operaciones lógicas de backend de forma privada y segura. \\
Server Action & Característica moderna de Next.js que procesa datos de formularios en el servidor directamente, evitando exponer código o endpoints al cliente. \\
RAG (Retrieval-Augmented Generation) & Técnica que alimenta al agente de IA con base documental interna de FABRIC antes de responder, garantizando respuestas precisas sobre doctrina y Oracle. \\
Embeddings & Traducción matemática de fragmentos de texto en vectores numéricos, permitiendo que la IA compare y encuentre textos similares de forma semántica. \\
PgVector / Pinecone & Motores especializados de base de datos vectoriales diseñados para almacenar embeddings y realizar búsquedas de contexto para la IA en milisegundos. \\
Fallback (Plan de respaldo) & Mecanismo que redirige una solicitud a un proveedor de IA secundario si el principal (Claude) sufre caídas de servicio o latencia severa. \\
Kill Switch & Clasificador de seguridad que detiene el procesamiento de la IA de inmediato si el usuario intenta realizar un "jailbreak" o pregunta sobre temas fuera de alcance. \\
High-Intent Routing & Análisis automático de la conversación de IA. Si detecta urgencia o problemas críticos en el ERP del usuario, activa botones directos para agendar llamadas. \\
Honeypot & Campo invisible para humanos que solo llenan los robots de spam. Si contiene datos, el sistema rechaza y bloquea el envío de forma silenciosa. \\
Rate limiting & Límite automático de peticiones por segundo por dirección IP para evitar ataques de denegación de servicio (DDoS) o spam masivo en los formularios. \\
Turnstile & Solución de CAPTCHA invisible desarrollada por Cloudflare. Bloquea bots sin molestar al usuario con imágenes confusas. \\
Lighthouse & Herramienta oficial de Google que audita y califica el rendimiento, accesibilidad, SEO y mejores prácticas técnicas del sitio web en una escala de 0 a 100. \\
LCP (Largest Contentful Paint) & Métrica de Lighthouse que mide la velocidad de carga visual. Evalúa el tiempo en que aparece el bloque principal de contenido en la pantalla. \\
CLS (Cumulative Layout Shift) & Métrica de Lighthouse que evalúa la estabilidad visual. Mide si los elementos gráficos cambian de posición de forma brusca o molesta durante la carga. \\
TTFB (Time to First Byte) & Tiempo que transcurre desde que el usuario hace clic hasta que el servidor envía el primer byte de información. Es una métrica de velocidad de respuesta del servidor. \\
FOUT (Flash of Unstyled Text) & Parpadeo visual molesto que ocurre cuando el navegador muestra una fuente genérica temporal antes de cargar la tipografía de diseño del sitio. \\
Webhook & Notificación automática en tiempo real enviada desde un sistema a otro ante un evento específico (ej. Calendly avisa a MongoDB que un lead agendó una cita). \\
Variables de entorno & Contenedor seguro en el servidor donde se guardan claves secretas, contraseñas de bases de datos y tokens de API, manteniéndolas fuera del código fuente. \\
Middleware & Capa intermedia que intercepta solicitudes de navegación (ej. verifica si el usuario tiene sesión activa antes de dejarlo ingresar al panel \texttt{/admin}). \\
\bottomrule
\end{longtable}


\section{Motor IA — especificación técnica}

El agente IA funciona como demostración técnica del dominio Oracle de FABRIC. Su diseño prioriza alcance controlado, trazabilidad y continuidad ante fallos de proveedor.

\subsection{Arquitectura del agente}

El flujo operativo de procesamiento de consultas del Agente Oracle sigue una secuencia estructurada de control para garantizar respuestas precisas, seguras y de baja latencia:

\begin{enumerate}
\item \textbf{Ingreso de Consulta:} El usuario escribe una pregunta técnica sobre Oracle o la firma en la interfaz de la terminal.
\item \textbf{Filtro de Alcance (Scope Filter):} Se intercepta la consulta del lado del servidor para verificar si está relacionada con el dominio técnico (Oracle EBS, Fusion Cloud, JDE, PeopleSoft, OCI) o institucional (FABRIC).
\item \textbf{Interruptor de Seguridad (Kill Switch):}
  \begin{itemize}
  \item Si la consulta está \textbf{fuera de alcance} o se detecta un intento de jailbreak, se detiene la ejecución inmediatamente, sugiriendo una vía de comunicación humana.
  \item Si la consulta es \textbf{válida}, se permite su avance al pipeline de datos.
  \end{itemize}
\item \textbf{Recuperación Aumentada por Generación (RAG):} Se consulta la base de datos vectorial para extraer los embeddings semánticos más relevantes de la documentación de Oracle y del material metodológico interno de FABRIC.
\item \textbf{Orquestación del Modelo (Pipeline de IA):}
  \begin{itemize}
  \item \textbf{Canal Primario:} Se envía la consulta junto al contexto del RAG a \textbf{Claude Sonnet}.
  \item \textbf{Primer Fallback:} Si el servicio primario reporta un error o excede un umbral de latencia de 3 segundos, se redirige la petición a \textbf{OpenAI GPT-4o}.
  \item \textbf{Segundo Fallback:} Si el fallo persiste, se emplea \textbf{Grok} como última instancia de recuperación.
  \end{itemize}
\item \textbf{Respuesta en Streaming:} Los datos generados se transmiten progresivamente al navegador mediante el hook \texttt{useChat} de Vercel AI SDK para reducir el tiempo percibido de carga.
\item \textbf{Enrutamiento de Alta Intención (High-Intent Routing):} Si el analizador semántico identifica que la consulta denota una urgencia crítica en el ERP del cliente, se activa y despliega de inmediato un llamado de acción directo para programar Office Hours.
\end{enumerate}

\subsection{Componentes técnicos}

\begin{longtable}{p{0.46\linewidth}p{0.46\linewidth}}
\toprule
\textbf{Componente} & \textbf{Implementación} \\
\midrule
Vercel AI SDK & Integración del hook \texttt{useChat} en el cliente para gestionar respuestas en streaming y controlar los estados de carga y error en la interfaz. \\
/api/ai & Endpoint en el servidor que implementa validación con Zod, filtrado de intención y streaming seguro de datos. \\
Orquestador multi-modelo & Módulo en \texttt{lib/ai.ts} encargado de enrutar consultas hacia Claude Sonnet y gestionar de forma reactiva los fallbacks hacia OpenAI o Grok. \\
RAG — contenido Oracle & Indexación y consulta de embeddings semánticos en base vectorial (PgVector en Supabase o Pinecone) con material técnico y doctrina. \\
Kill Switch & Clasificador semántico server-side que intercepta y bloquea consultas inapropiadas o ajenas al dominio temático del negocio. \\
High-Intent Routing & Filtro de palabras clave y análisis de contexto que dispara dinámicamente un llamado de acción prioritario para agendar reuniones de consultoría. \\
Terminal UI & Componente React de consola que emula una terminal de comandos, con efectos de escritura (typewriter), indicador parpadeante y preselección de escenarios. \\
\bottomrule
\end{longtable}

\subsection{Contenido RAG pendiente de carga}

La arquitectura RAG queda preparada para recibir material validado. La carga de vectores se realiza cuando los documentos estén disponibles, sin modificar el flujo principal del agente.

\begin{longtable}{p{0.46\linewidth}p{0.46\linewidth}}
\toprule
\textbf{Categoría de contenido} & \textbf{Estado / fuente} \\
\midrule
Documentación Oracle Fusion Cloud: Financials, SCM, HCM y Procurement & Pendiente · Oracle docs públicos + material interno \\
Casos FABRIC: APE Plazas y Aplazo & Pendiente · Julio entrega material técnico \\
Doctrina FABRIC: cinco compromisos operativos & Pendiente · documento de dirección \\
Papers de investigación & Pendiente · PDFs internos \\
Patrones de fracaso Oracle en LATAM & Pendiente · material interno \\
FSOs: seis Fabric Solution Objects & Pendiente · confirmación de contenido \\
\bottomrule
\end{longtable}

Proceso de carga: /scripts/index-rag.ts toma los documentos, genera embeddings mediante OpenAI API y almacena los vectores en PgVector. El script se ejecuta al cargar nuevo contenido o al actualizar el corpus.

\section{Panel de administración — especificación técnica}

El panel admin permite operar leads, métricas, slots y capacidad desde una interfaz privada, sin tocar código ni redeployar el sitio.

\subsection{Arquitectura y acceso}

\begin{longtable}{p{0.46\linewidth}p{0.46\linewidth}}
\toprule
\textbf{Elemento} & \textbf{Especificación} \\
\midrule
URL de acceso & fabricsoft.com.mx/admin, sin enlace visible en el sitio público. \\
Autenticación & Clerk con restricción a emails @fabricsoft.com.mx y 2FA. \\
Middleware & middleware.ts intercepta /admin/* y valida sesión antes de renderizar. \\
Roles & Super Admin con acceso total; Admin con permisos limitados. \\
Log de acciones & Registro inmutable en MongoDB con usuario, acción, fecha y cambios. \\
Propagación & Cambios de métricas visibles mediante ISR en máximo una hora, sin redeploy. \\
\bottomrule
\end{longtable}

\subsection{Funcionalidades completas}

\subsubsection{Gestión de leads}

\begin{enumerate}
\item Vista tabular con nombre, empresa, cargo, revenue declarado, tipo de iniciativa, score y fecha.
\item Filtros por tipo, score, industria y fecha.
\item Acciones: aprobar, rechazar y ver detalle completo.
\item Indicador visual de calidad: champagne para alta prioridad, gris para media y rojo para rechazo recomendado.
\end{enumerate}

\subsubsection{Edición de métricas en vivo}

\begin{enumerate}
\item Rescue Counter: rescates, horas ahorradas, reportes eliminados y cierres estabilizados.
\item Wait List: proyectos activos, conteo de lista y próxima ventana de admisión.
\item Office Hours: crear/eliminar slots, editar URL de Calendly y marcar disponibilidad.
\item Cambios guardados en Sanity y reflejados sin deploy.
\end{enumerate}

\subsubsection{Control de capacidad operativa}

\begin{enumerate}
\item Variable de máximo de proyectos configurable desde la UI.
\item Dashboard de estado con slots disponibles, wait list y proyectos activos.
\item Toggle manual para activar Sold Out o forzar modo Wait List.
\end{enumerate}

\subsection{Wireframe del dashboard admin}

El centro de control administrativo adopta un diseño de panel dividido (Sidebar + Main Content Area) optimizado para la gestión operativa diaria:

\begin{description}
\item[Barra Lateral de Navegación (Sidebar Izquierda Fija):] \ 
  \begin{itemize}
  \item Identificador de Marca: Logo corporativo de FABRIC.
  \item Módulos de Acceso: Enlaces directos a \textbf{Dashboard}, \textbf{Leads}, \textbf{Métricas}, \textbf{Capacidad} y \textbf{Logs de Auditoría}.
  \end{itemize}

\item[Vista 1: Dashboard Resumen (Área Principal por Defecto):] \ 
  \begin{itemize}
  \item Mapeo de Métricas: Fichas de estado con leads registrados hoy y solicitudes pendientes de revisión.
  \item Estado de Operación: Indicador visual de proyectos activos frente al límite configurado y slots de Office Hours disponibles.
  \item Actividad Reciente: Tabla compacta con los últimos 5 leads y botones rápidos de aprobación.
  \end{itemize}

\item[Vista 2: Gestión de Leads (Tabla Detallada):] \ 
  \begin{itemize}
  \item Panel de Filtros: Selectores de búsqueda por tipo de iniciativa, score de prioridad, industria y rango de fechas.
  \item Tabla de Información: Lista detallada que muestra Nombre, Empresa, Cargo, Score, Fecha y opciones para inspeccionar respuestas completas de los wizards.
  \end{itemize}

\item[Vista 3: Editor de Métricas en Vivo:] \ 
  \begin{itemize}
  \item Formularios de Entrada: Campos numéricos para modificar directamente las cifras del Rescue Counter (horas ahorradas, cierres estabilizados) y parámetros de la lista de espera.
  \item Administrador de Slots: Tabla interactiva para dar de alta o baja enlaces de Calendly en las Office Hours del mes.
  \end{itemize}

\item[Vista 4: Configuración de Capacidad y Estado:] \ 
  \begin{itemize}
  \item Control Deslizante (Slider): Establecer la capacidad máxima de proyectos simultáneos.
  \item Toggle de Forzado Manual: Activación manual e inmediata de la Wait List o estado de "Sold Out".
  \item Monitor del Estado del Sitio: Visualización en tiempo real del estado de cara al público (Normal, Wait List, Sold Out).
  \end{itemize}
\end{description}

\section{Integración y consumo de datos}

\subsection{Servicios y costos}

\begin{longtable}{p{0.31\linewidth}p{0.31\linewidth}p{0.31\linewidth}}
\toprule
\textbf{Servicio} & \textbf{Rol en V1} & \textbf{Costo mes 1} \\
\midrule
Sanity CMS & Contenido editable: casos, papers, FSOs, métricas y slots & Gratuito \\
MongoDB Atlas & Leads, diagnósticos, resultados TCO, bookings, descargas y logs & Gratuito 512 MB \\
Resend & Templates de email transaccional & Gratuito 3,000/mes \\
HubSpot CRM & Registro de leads y estados comerciales & Gratuito \\
Calendly & Booking de Office Hours & Premium, costo por confirmar \\
Anthropic API & Modelo primario del agente IA & Variable por uso \\
OpenAI API & Fallback y embeddings RAG & Variable por uso \\
Grok API & Segundo fallback del agente IA & Variable por uso \\
Clerk & Autenticación del panel admin & Hobby tier \\
Slack Webhook & Notificación interna de nuevos leads & Gratuito \\
Cloudflare Turnstile & CAPTCHA invisible & Gratuito \\
Plausible Analytics & Analítica privacy-first y eventos custom & USD 9/mes \\
Sentry & Monitoreo de errores & Gratuito \\
Total fijo estimado & Depende del uso del agente IA & USD 9/mes + APIs IA \\
\bottomrule
\end{longtable}

\subsection{Modelo documental MongoDB Atlas}

MongoDB Atlas concentra los datos operativos y transaccionales del ecosistema FABRIC. Sanity conserva el contenido editorial; MongoDB conserva aquello que cambia por interacción del usuario, reglas de admisión, operación comercial, agenda, auditoría y métricas vivas.

\diagramimage{modelo-bd-mongodb.png}{Modelo lógico documental para MongoDB Atlas.}{fig:modelo-bd}

\subsubsection{Colecciones propuestas}

\begin{longtable}{p{0.18\linewidth}p{0.30\linewidth}p{0.24\linewidth}p{0.18\linewidth}}
\toprule
\textbf{Colección} & \textbf{Propósito} & \textbf{Campos clave} & \textbf{Índices recomendados} \\
\midrule
leads & Registro maestro del prospecto calificado o rechazado silenciosamente. & contact, company, role, revenueBand, source, score, status, tags, createdAt. & unique email, company.domain, status + score, createdAt. \\
diagnosticRuns & Resultados del Rescue Assessment y señales de dolor operativo. & leadId, answers, painSignals, score, recommendation, submittedAt. & leadId, score, submittedAt. \\
tcoComparisons & Entradas y resultados de comparadores ERP/OCI. & leadId, currentStack, annualCost, userCount, savingsEstimate, breakevenMonths. & leadId, createdAt, savingsEstimate. \\
aiConversations & Conversaciones con el Agente IA Oracle y señales high-intent. & leadId, sessionId, messages, scopeStatus, provider, tokenUsage, highIntent, createdAt. & sessionId, leadId, highIntent, createdAt. \\
paperRequests & Solicitudes de papers y materiales técnicos. & leadId, paperSlug, email, company, consent, deliveredAt. & email + paperSlug, leadId, deliveredAt. \\
officeHourBookings & Reservas o intentos de reserva de Office Hours. & leadId, calendlyEventId, slot, status, qualificationSnapshot. & calendlyEventId, leadId, slot.start, status. \\
capacityState & Estado operativo editable desde admin. & activeProjects, maxProjects, waitListCount, officeHourSlots, forcedMode, updatedAt. & forcedMode, updatedAt. \\
publicMetrics & Métricas públicas mostradas en la home y el portal de transparencia. & rescueCount, hoursSaved, reportsEliminated, stabilizedClosings, fixedPricePercentage, npsScore, retentionRate, seniorConsultantPercentage, averageResponseTime, source, updatedAt. & updatedAt. \\
adminAuditLogs & Bitácora inmutable de acciones administrativas. & adminUserId, action, targetCollection, targetId, before, after, createdAt. & adminUserId, targetId, createdAt. \\
\bottomrule
\end{longtable}

\subsubsection{Documento base: lead}

\begin{Verbatim}[fontsize=\small, breaklines=true]
{
  "_id": "ObjectId",
  "contact": {
    "name": "string",
    "email": "string",
    "phone": "string",
    "role": "CFO | CTO | CIO | Founder | Other"
  },
  "company": {
    "name": "string",
    "domain": "string",
    "industry": "financial | real-estate | logistics | other",
    "revenueBand": "50M-100M | 100M-500M | 500M+ | unknown",
    "country": "MX | LATAM | other"
  },
  "qualification": {
    "score": 0,
    "status": "new | qualified | rejected | waitlist | booked",
    "initiative": "rescue | migration | optimization | advisory",
    "initiativeTimeline": "under_6_months | over_6_months",
    "execSponsorship": "CFO_AND_CTO | OTHER",
    "isCorporateEmail": true,
    "rejectionReason": "string | null"
  },
  "source": {
    "entryPoint": "/diagnostico",
    "utm": {},
    "firstSeenAt": "ISODate"
  },
  "integrations": {
    "hubspotContactId": "string",
    "slackNotificationId": "string"
  },
  "createdAt": "ISODate",
  "updatedAt": "ISODate"
}
\end{Verbatim}

\subsubsection{Justificación técnica y criterios de diseño de la base de datos}

El ecosistema de FABRIC emplea \textbf{MongoDB Atlas} como su base de datos operativa debido a las siguientes razones de ingeniería:
\begin{itemize}
\item \textbf{Flexibilidad de Esquema:} Las interacciones del Agente IA (conversaciones completas en JSON), las respuestas dinámicas del Rescue Assessment (que varían según el dolor operativo) y los metadatos UTM son semiestructurados. MongoDB permite almacenar esta información de forma nativa sin necesidad de migraciones complejas de esquemas SQL.
\item \textbf{Desacoplamiento de Lectura/Escritura:} La estructura documental permite almacenar snapshots embebidos del lead en las colecciones de eventos (diagnósticos, bookings, etc.), evitando consultas pesadas con múltiples uniones (JOINs) y optimizando el tiempo de respuesta en el Edge.
\item \textbf{Criterios Específicos por Colección:}
  \begin{itemize}
  \item \texttt{leads}: Actúa como el \emph{Single Source of Truth} del prospecto, centralizando su perfil corporativo e historial de calificación para evitar silos de datos.
  \item \texttt{diagnosticRuns}, \texttt{tcoComparisons} y \texttt{aiConversations}: Se separan del documento base del lead para evitar el crecimiento ilimitado de documentos únicos y facilitar análisis aislados de conversión e indexación RAG.
  \item \texttt{capacityState} y \texttt{publicMetrics}: Están diseñadas para almacenamiento de baja latencia. \texttt{publicMetrics} reduce la carga en la BD operativa al consolidar las métricas que consume la página pública de transparencia.
  \item \texttt{adminAuditLogs}: Es estrictamente de tipo \emph{append-only} para garantizar la inmutabilidad y la trazabilidad de cualquier alteración manual de leads o slots por parte del staff.
  \end{itemize}
\end{itemize}

Las Server Actions y las API Routes validan cada payload con esquemas de Zod y sanitizan las entradas antes de interactuar con el motor de base de datos para prevenir inyecciones.

\subsubsection{Flujo de trabajo de admisión y ciclo de vida del lead (Customer Journey)}

El motor de admisión funciona como un \emph{gatekeeper} server-side que controla el flujo de contactos desde el primer clic hasta la asignación de recursos de ingeniería. El ciclo opera bajo la siguiente secuencia lógica:

\begin{enumerate}
\item \textbf{Entrada y Captura (First Touch):} El usuario interactúa con la plataforma a través de los wizard de diagnóstico (/diagnostico), comparador (/comparator/erp), solicitud de papers (/investigacion) o el Agente IA.
\item \textbf{Filtro de Admisión Server-Side (Gating):} Al enviar los datos, Next.js procesa la petición de forma segura:
  \begin{itemize}
  \item Valida que el correo electrónico no pertenezca a dominios públicos o de spam.
  \item Verifica los ingresos corporativos declarados (debe ser $\ge$ 50M USD).
  \item Comprueba las verticales aprobadas (Fintech/Servicios Financieros, Inmobiliario, Logística).
  \item Confirma el patrocinio ejecutivo conjunto de la iniciativa (CFO + CTO).
  \end{itemize}
\item \textbf{Bifurcación de Calificación:}
  \begin{itemize}
    \item \textbf{Lead Calificado (Estatus \texttt{qualified}):} Si cumple con todos los filtros, el sistema habilita dinámicamente la redirección al widget de Calendly para agendar sus \emph{Office Hours} (el estado pasa a \texttt{booked} al reservar).
    \item \textbf{Lead No Calificado (Estatus \texttt{waitlist} o \texttt{rejected}):} Si el prospecto no cumple los criterios (e.g. es una empresa de menor facturación o fuera de verticales), el sistema bloquea silenciosamente el agendamiento y lo redirige a la lista de espera (\texttt{waitlist}), enviándole recursos automáticos de apoyo por correo.
  \end{itemize}
\item \textbf{Sincronización de Sistemas y Notificaciones:}
  \begin{itemize}
  \item Se registra el lead y su evento relacionado en MongoDB Atlas.
  \item Se invoca la API de HubSpot CRM para crear o actualizar el contacto con sus etiquetas de calificación.
  \item Se detona un webhook para alertar al canal interno de Slack de FABRIC en tiempo real.
  \item Se utiliza la API de \textbf{Resend} para despachar de inmediato el material solicitado (PDF de paper o reporte de diagnóstico).
  \end{itemize}
\end{enumerate}

\subsection{Flujos de lead magnets}

\subsubsection{Flujo 1 — Rescue Assessment (/diagnostico)}
El embudo de autodiagnóstico de fallas en implementaciones de Oracle Fusion opera en base al siguiente flujo interactivo y transaccional:
\begin{enumerate}
\item \textbf{Cuestionario:} El usuario responde un wizard de 12 preguntas cerradas sobre sus problemas de sistema, cierre mensual y patrocinio.
\item \textbf{Evaluación Reactiva:} El navegador computa un score técnico preliminar y asigna un nivel de prioridad al caso.
\item \textbf{Validación y Guardado:} Al enviar el formulario con datos de contacto, Next.js valida el correo corporativo y la integridad del payload con Zod antes de guardarlo en MongoDB Atlas.
\item \textbf{Notificación e Integración:} Se crea un perfil calificado en HubSpot CRM, se activa un webhook para notificar al canal de Slack de FABRIC y se despacha un correo de confirmación mediante Resend.
\end{enumerate}

\subsubsection{Flujo 2 — ERP TCO Comparator (/comparator/erp)}
El estimador de costo de propiedad (TCO) permite a directivos evaluar el beneficio financiero de migrar a Oracle Fusion:
\begin{enumerate}
\item \textbf{Wizard de Entrada:} Formulario secuencial de 8 pasos que captura el ERP actual, número de usuarios e inversiones en licenciamiento, soporte e infraestructura.
\item \textbf{Cálculo en Tiempo Real:} A través de un módulo especializado, se procesan los ahorros potenciales basados en benchmarks sectoriales de ahorro y retorno de inversión.
\item \textbf{Gating de Datos:} Los resultados finales se muestran tras la captura de un correo corporativo válido.
\item \textbf{Distribución de Resultados:} Se persisten las variables de la simulación en la base de datos, se actualiza el CRM y se remite un PDF con el análisis comparativo formal al usuario vía correo electrónico.
\end{enumerate}

\subsubsection{Flujo 3 — Wait List (/aplicar)}
La lista de espera filtra prospectos cuando la capacidad de ingeniería está colmada o el perfil corporativo requiere validación manual:
\begin{enumerate}
\item \textbf{Formulario Unificado:} El usuario ingresa información detallada en un formulario estructurado de 10 campos críticos.
\item \textbf{Acción del Servidor:} Una Server Action en Next.js valida la información y procesa el honeypot para mitigar spam de forma segura.
\item \textbf{Almacenamiento:} Se guarda el lead con el estado de "lista de espera" en MongoDB Atlas y se sincroniza con la lista homónima en HubSpot CRM.
\item \textbf{Feedback al Usuario:} Se actualiza la interfaz mostrando la ventana estimada de admisión y se le remite una notificación formal.
\end{enumerate}

\subsubsection{Flujo 4 — Office Hours (/office-hours)}
El canal directo de agendamiento para empresas que requieren remediación de alta prioridad:
\begin{enumerate}
\item \textbf{Filtro de Admisión:} El sistema ejecuta un cuestionario de precalificación validando cargo (C-Level) y revenue corporativo (USD 50M+).
\item \textbf{Bifurcación del Flujo:} 
  \begin{itemize}
  \item Si el usuario cumple con los umbrales de negocio, el sistema habilita y renderiza dinámicamente el widget inline de Calendly.
  \item Si la información no es calificada, se realiza una redirección hacia el flujo de la Wait List para capturar el contacto sin interrumpir el proceso.
  \end{itemize}
\item \textbf{Sincronización:} Tras reservar la fecha, un webhook de Calendly notifica a la base de datos para registrar el booking y actualizar el estado del lead a "booked".
\end{enumerate}

\subsubsection{Flujo 5 — Papers Técnicos (/investigacion)}
La distribución de materiales de investigación técnica especializada está diseñada para consolidar la autoridad de la firma ante directivos:
\begin{enumerate}
\item \textbf{Catálogo de Publicaciones:} La plataforma expone tres papers clave:
  \begin{itemize}
  \item \emph{Paper 01: Por qué fallan los go-live de Oracle Fusion} (Análisis empírico de 47 casos de remediación en LATAM).
  \item \emph{Paper 02: IA aplicada al cierre contable en Fusion Cloud} (Framework de automatización y casos de uso prácticos).
  \item \emph{Paper 03: Modelo de entrega en primer ciclo crítico} (La doctrina operativa de FABRIC contractualizada).
  \end{itemize}
\item \textbf{Control de Acceso Riguroso (Gating):} La descarga directa en el navegador está completamente deshabilitada en el cliente. El usuario debe enviar un formulario validado mediante una Server Action.
\item \textbf{Validación del Lado del Servidor:} Zod y el motor de precalificación verifican que el correo sea de dominio corporativo, el cargo corresponda a niveles directivos y la empresa pertenezca a las verticales aprobadas.
\item \textbf{Despacho Asíncrono de Información:} Si la validación es exitosa, se registra en la colección \texttt{paperRequests} y se invoca la API de \textbf{Resend} para despachar el PDF al correo del destinatario. En caso de no calificar, el lead se guarda con estatus \texttt{waitlist} y se le notifica la restricción de envío de forma diplomática.
\end{enumerate}

\subsubsection{Flujo 6 — Agente IA (S04)}
El asistente conversacional técnico sirve como primer punto de validación autónoma de la doctrina y pericia Oracle de la firma:
\begin{enumerate}
\item \textbf{Validación de Consulta:} Se clasifica la intención de la pregunta del usuario en el endpoint del servidor.
\item \textbf{Recuperación de Contexto:} Se ejecuta una búsqueda semántica de documentos relacionados en la base vectorial (RAG).
\item \textbf{Generación de Respuesta:} El orquestador de modelos genera la respuesta y la transmite en streaming utilizando el hook \texttt{useChat} del SDK.
\item \textbf{Escalación Comercial:} Si la conversación revela señales críticas del negocio (e.g. "cierres inestables", "retrasos de producción"), el enrutador habilita de forma visual un botón directo hacia las Office Hours.
\end{enumerate}

\subsubsection{Flujo 7 — Portal de Transparencia (/transparencia)}
El portal de transparencia expone públicamente los indicadores de éxito y auditoría operativa de FABRIC:
\begin{enumerate}
\item \textbf{Métricas Clave en Vivo:} Visualización de indicadores reales de desempeño:
  \begin{itemize}
  \item \emph{Proyectos entregados en primer ciclo crítico (100\%):} Evidencia empírica de cumplimiento del principio 1.
  \item \emph{Proyectos bajo esquema Fixed-Price sin desvíos de costo (100\%):} Reflejo del cumplimiento del principio 3.
  \item \emph{NPS Promedio de clientes activos (mínimo 90 pts):} Calidad del servicio.
  \item \emph{Porcentaje de personal con seniority mayor a 8 años (100\%):} Cumplimiento del principio 2.
  \item \emph{Tiempo medio de respuesta operativa a incidencias críticas.}
  \end{itemize}
\item \textbf{Estrategia de Actualización (ISR):} Los datos son administrados de forma privada por el equipo desde el panel de administración (\texttt{/admin}) o Sanity CMS. La página pública \texttt{/transparencia} se regenera bajo demanda en el Edge (ISR con revalidación de 3600 segundos), garantizando velocidad instantánea de carga y nulo impacto en el rendimiento de la base de datos operativa.
\end{enumerate}

\section{Infraestructura, rendimiento y despliegue}

\subsection{Hosting y entornos}

\begin{enumerate}
\item \textbf{Servicio de Alojamiento:} Implementación sobre la plataforma Vercel Pro, aprovechando Edge Functions, optimizaciones integradas de caché para Incremental Static Regeneration (ISR) y generación automática de URLs de previsualización por rama Git.
\item \textbf{Entorno de Preproducción:} Configuración del subdominio de pruebas \texttt{equipoa.fabricsoft.com.mx} para revisiones y pruebas del equipo interno antes del paso a producción.
\item \textbf{Seguridad y Conectividad DNS:} Delegación de zonas DNS a Cloudflare para optimizar la latencia mediante HTTP/3, activar reglas de firewall perimetral y administrar certificados SSL automáticos.
\item \textbf{Aislamiento de Ambientes:} Separación de claves y configuraciones críticas en tres variables independientes (Desarrollo, Staging y Producción).
\end{enumerate}

\subsection{Estrategia Lighthouse 95+}

\begin{enumerate}
\item \textbf{Gestión de Fuentes:} Uso del módulo nativo \texttt{next/font} para cargar tipografías en tiempo de compilación y servirlas de forma local, eliminando el parpadeo visual (FOUT) y previniendo caídas en la métrica CLS.
\item \textbf{Optimización de Recursos Gráficos:} Uso exclusivo del componente \texttt{next/image} de React para forzar la conversión automática de archivos a WebP, escalado responsivo inteligente e inicialización diferida (lazy loading).
\item \textbf{Renderizado Estático Inteligente:} Priorización de Static Generation (SSG) en páginas estáticas del sitio e ISR en aquellas secciones que consumen variables modificables de Sanity o MongoDB Atlas.
\item \textbf{Reducción del Tamaño de Código (Bundles):} Exclusión de frameworks pesados de UI, purgado estricto de clases CSS no utilizadas mediante Tailwind compiler y optimización de animaciones de Framer Motion.
\item \textbf{Asincronismo de Datos:} Ejecución distribuida y streaming asíncrono para las respuestas de la terminal de IA para impedir la saturación del hilo principal del navegador.
\end{enumerate}

\subsection{Seguridad}

\begin{enumerate}
\item HTTP Headers: CSP, HSTS, X-Frame-Options y X-Content-Type-Options.
\item Clerk middleware protege /admin/* y rutas API administrativas.
\item Turnstile CAPTCHA invisible en formularios públicos.
\item API keys de IA, Clerk y servicios externos siempre en variables de entorno server-side.
\item Cumplimiento LFPDPPP: aviso de privacidad en footer y consentimiento explícito en formularios.
\end{enumerate}

\subsection{SEO técnico}

\begin{enumerate}
\item Meta tags únicos por página: title, description y canonical URL.
\item Open Graph con imagen 1200x630 px, fondo \#0A0A0A, logo FABRIC y acento champagne.
\item sitemap.xml generado automáticamente por Next.js en build.
\item robots.txt indexa páginas públicas y excluye /api/*, /admin/* y /studio.
\item Schema.org ProfessionalService con founder, areaServed y knowsAbout Oracle.
\item Lighthouse SEO 100 verificado antes de cada demo.
\end{enumerate}

\end{document}
